\chapter{Resultados de Análisis Eléctricos}

\section{Flujos de Carga en Operación Normal}

Los análisis de flujo de carga se realizaron en condición normal de operación, con el fin de evaluar el desempeño del sistema ante la entrada en operación del proyecto fotovoltaico.

\subsection{Análisis de Tensiones}

La entrada en operación del proyecto ocasiona variaciones de tensión en los nodos cercanos al punto de conexión, con un cambio máximo de 0,00062 p.u., valor que no genera ninguna violación a los límites operativos normales del sistema establecidos en el Código de Redes.

\begin{table}[H]
    \centering
    \caption{Variación de tensión a las 9 a.m.}
    \begin{tabular}{lccc}
    \toprule
    \textbf{Nodo} & \textbf{Sin SSFV} & \textbf{Con SSFV} & \textbf{Variación} \\
    \textbf{} & \textbf{(p.u.)} & \textbf{(p.u.)} & \textbf{(p.u.)} \\
    \midrule
    3272869 & 0.99470 & 0.99532 & 0.00062 \\
    3272966 & 0.99470 & 0.99532 & 0.00062 \\
    3272761 & 0.99470 & 0.99532 & 0.00061 \\
    3272664 & 0.99471 & 0.99532 & 0.00061 \\
    3272567 & 0.99471 & 0.99532 & 0.00060 \\
    1065971 & 0.99471 & 0.99531 & 0.00060 \\
    2479567 & 0.99471 & 0.99531 & 0.00060 \\
    \bottomrule
    \end{tabular}
\end{table}

\begin{table}[H]
    \centering
    \caption{Variación de tensión a las 12 p.m.}
    \begin{tabular}{lccc}
    \toprule
    \textbf{Nodo} & \textbf{Sin SSFV} & \textbf{Con SSFV} & \textbf{Variación} \\
    \textbf{} & \textbf{(p.u.)} & \textbf{(p.u.)} & \textbf{(p.u.)} \\
    \midrule
    3272869 & 0.99432 & 0.99494 & 0.00062 \\
    3272966 & 0.99432 & 0.99494 & 0.00062 \\
    3272761 & 0.99432 & 0.99494 & 0.00061 \\
    3272664 & 0.99433 & 0.99494 & 0.00061 \\
    3272567 & 0.99433 & 0.99494 & 0.00060 \\
    1065971 & 0.99433 & 0.99493 & 0.00060 \\
    2479567 & 0.99433 & 0.99493 & 0.00060 \\
    \bottomrule
    \end{tabular}
\end{table}

\begin{table}[H]
    \centering
    \caption{Variación de tensión a las 3 p.m.}
    \begin{tabular}{lccc}
    \toprule
    \textbf{Nodo} & \textbf{Sin SSFV} & \textbf{Con SSFV} & \textbf{Variación} \\
    \textbf{} & \textbf{(p.u.)} & \textbf{(p.u.)} & \textbf{(p.u.)} \\
    \midrule
    3272869 & 0.99392 & 0.99454 & 0.00061 \\
    3272966 & 0.99392 & 0.99454 & 0.00062 \\
    3272761 & 0.99393 & 0.99454 & 0.00061 \\
    3272664 & 0.99393 & 0.99454 & 0.00061 \\
    3272567 & 0.99393 & 0.99454 & 0.00060 \\
    1065971 & 0.99393 & 0.99453 & 0.00060 \\
    2479567 & 0.99393 & 0.99453 & 0.00060 \\
    \bottomrule
    \end{tabular}
\end{table}

\textbf{Conclusión:} Las variaciones de tensión están dentro de los límites permitidos (0.90 - 1.10 p.u.).

\subsection{Cargabilidad de Líneas de Distribución}

Se evaluaron las líneas cercanas al proyecto para identificar si existe variación en la dirección de los flujos de potencia y reducciones o aumentos en la cargabilidad de las líneas. Se identificó una variación máxima de 1.725\% ante la entrada del proyecto.

\textbf{Observación importante:} No existe un aumento de la cargabilidad en ninguna línea. Esto indica que toda la energía generada es consumida de manera local, reduciendo de esta manera la cargabilidad y por ende las pérdidas en las líneas.

\begin{table}[H]
    \centering
    \caption{Cargabilidad de líneas a las 12 p.m. (máxima variación)}
    \begin{tabular}{lccc}
    \toprule
    \textbf{Línea} & \textbf{Sin SSFV (\%)} & \textbf{Con SSFV (\%)} & \textbf{Variación (\%)} \\
    \midrule
    5291 & 4.565 & 2.839 & -1.725 \\
    5292 & 4.565 & 2.840 & -1.725 \\
    5294 & 4.565 & 2.840 & -1.725 \\
    5295 & 4.566 & 2.841 & -1.725 \\
    5296 & 4.566 & 2.841 & -1.724 \\
    26921 & 2.812 & 1.303 & -1.510 \\
    5373 & 4.566 & 2.842 & -1.724 \\
    \bottomrule
    \end{tabular}
\end{table}

\textbf{Conclusión:} En ningún caso se violan los límites de cargabilidad del sistema de distribución de ESSA.

\subsection{Cargabilidad de Transformadores}

Se cuenta con un transformador principal de 13.2 kV/220V en la cabecera de circuito que suministra energía en media tensión. Se observa un aumento máximo de la cargabilidad del transformador de 28.833\%.

\begin{table}[H]
    \centering
    \caption{Cargabilidad del transformador de conexión}
    \begin{tabular}{lccc}
    \toprule
    \textbf{Hora} & \textbf{Sin SSFV (\%)} & \textbf{Con SSFV (\%)} & \textbf{Variación (\%)} \\
    \midrule
    9:00 a.m. & 29.645 & 53.655 & 24.010 \\
    12:00 p.m. & 28.598 & 54.403 & 25.805 \\
    15:00 (3 p.m.) & 27.000 & 55.833 & 28.833 \\
    \bottomrule
    \end{tabular}
\end{table}

\textbf{Conclusión:} La cargabilidad siempre se mantiene por debajo de los límites operativos normales, aún más se encuentra por debajo del 60\%, lo que deja una capacidad extra que puede ser aprovechada para una generación mayor o para más demanda en la zona.

\section{Resultados del Estudio de Cortocircuito}

Se realizó el cálculo de los niveles de cortocircuito trifásico y monofásico del NODO 3272966 del OR ESSA considerando la entrada en servicio del proyecto de generación fotovoltaica de 120 kW.

\subsection{Cortocircuito Monofásico}

\begin{table}[H]
    \centering
    \caption{Cortocircuito monofásico a las 12 p.m. (nodos principales)}
    \begin{tabular}{lccc}
    \toprule
    \textbf{Nodo} & \textbf{Sin SSFV (kA)} & \textbf{Con SSFV (kA)} & \textbf{Variación (kA)} \\
    \midrule
    3272869 & 3.15235 & 3.15235 & 0.00000000 \\
    3272966 & 3.14984 & 3.14984 & 0.00000000 \\
    3272761 & 3.15571 & 3.15571 & 0.00000000 \\
    3272664 & 3.16218 & 3.16218 & 0.00000000 \\
    3272567 & 3.17766 & 3.17766 & 0.00000000 \\
    \bottomrule
    \end{tabular}
\end{table}

\subsection{Cortocircuito Trifásico}

\begin{table}[H]
    \centering
    \caption{Cortocircuito trifásico a las 12 p.m. (nodos principales)}
    \begin{tabular}{lccc}
    \toprule
    \textbf{Nodo} & \textbf{Sin SSFV (kA)} & \textbf{Con SSFV (kA)} & \textbf{Variación (kA)} \\
    \midrule
    3272869 & 5.12685 & 5.12685 & 0.00000000 \\
    3272966 & 5.12027 & 5.12027 & 0.00000000 \\
    3272761 & 5.13564 & 5.13564 & 0.00000000 \\
    3272664 & 5.15265 & 5.15265 & 0.00000000 \\
    3272567 & 5.19351 & 5.19351 & 0.00000000 \\
    \bottomrule
    \end{tabular}
\end{table}

\textbf{Análisis de Resultados:} A partir del análisis de las corrientes de cortocircuito se identificó una variación máxima de 0.00000003 kA, esto representa una variación de menos del 0.01\% respecto al caso base.

\textbf{Conclusión:} Las corrientes de cortocircuito no presentan un aumento significativo, debido a que la tecnología de generación empleada (inversores solares) limita la corriente de inyección hacia el sistema. No es necesario realizar modificaciones en las protecciones existentes.

\section{Protección en el Punto de Conexión}

\subsection{Cálculo de Corriente Nominal}

La corriente máxima entregada por el generador fotovoltaico a potencia nominal por el lado de baja del transformador se calcula como:

\begin{equation}
I_L = \frac{S}{\sqrt{3} \times V_l \times FP}
\end{equation}

Donde:
\begin{itemize}
    \item $I_L$: Corriente nominal
    \item $S$: Potencia aparente del generador = 120 kVA
    \item $V_l$: Tensión de línea en media tensión = 220 V
    \item $FP$: Factor de potencia = 1
\end{itemize}

\begin{equation}
I_L = \frac{120}{\sqrt{3} \times 0.22} = 393.65 \text{ A}
\end{equation}

\subsection{Requisitos de Protección}

Según el Acuerdo 1322 del Consejo Nacional de Operación (CNO), la conexión de generadores al sistema eléctrico debe cumplir con los siguientes requisitos de protecciones:

\begin{table}[H]
    \centering
    \caption{Requerimientos de protección por nivel de tensión}
    \begin{tabular}{lccc}
    \toprule
    \textbf{Función de Protección} & \textbf{Umbral} & \textbf{Tiempo} & \textbf{Unidad} \\
    \midrule
    Baja tensión (ANSI 27) - Etapa 1 & 0.85 & 2.0 & p.u. / s \\
    Baja tensión (ANSI 27) - Etapa 2 & 0.50 & 0.2 & p.u. / s \\
    Sobretensión (ANSI 59) - Etapa 1 & $\geq 1.15$ & 2.0 & p.u. / s \\
    Sobretensión (ANSI 59) - Etapa 2 & $\geq 1.20$ & 0.1-0.2 & p.u. / s \\
    Subfrecuencia (ANSI 81U) & 57 & 0.5 & Hz / s \\
    Sobrefrecuencia (ANSI 81O) & 63 & 0.5 & Hz / s \\
    \bottomrule
    \end{tabular}
\end{table}

\subsection{Validación de Protecciones}

\begin{table}[H]
    \centering
    \caption{Validación de requisitos de protección}
    \begin{tabular}{lc}
    \toprule
    \textbf{Requerimiento de Protección} & \textbf{Validación} \\
    \midrule
    Baja tensión (ANSI 27) & Cumple \\
    Sobrepotencia adelante (ANSI 32) & No Aplica \\
    Sobretensión (ANSI 59) & Cumple \\
    Frecuencia (ANSI 81 U/O) & Cumple \\
    Anti - Isla & Cumple \\
    \bottomrule
    \end{tabular}
\end{table}

\textbf{Conclusión:} No es necesario realizar modificaciones en las protecciones.

\subsection{Elemento de Interrupción}

Para el cálculo de protecciones contra sobrecorrientes en el sistema solar, la protección del sistema se compone de:

\begin{itemize}
    \item \textbf{Protección de inversores:} $3 \times (140-200)$ A
    \item \textbf{Protección de acometida:} $3 \times (350-500)$ A
\end{itemize}

Dicha protección es tomada como medio de seccionamiento y protección en punto de conexión. Los tiempos de respuesta del equipo de protección seleccionado cumplen con la norma IEC 60909 e IEEE 242.

\section{Cálculo de Pérdidas}

Se realizó el cálculo de pérdidas de potencia con y sin la presencia del proyecto solar fotovoltaico dispuesto a conectarse en el NODO 3272966 del OR ESSA.

\subsection{Resumen de Pérdidas Totales}

\begin{table}[H]
    \centering
    \caption{Pérdidas totales del sistema}
    \begin{tabular}{lccc}
    \toprule
    \textbf{Hora} & \textbf{Con SSFV (kW)} & \textbf{Sin SSFV (kW)} & \textbf{Reducción (kW)} \\
    \midrule
    9:00 a.m. & 4.47 & 4.77 & 0.30 \\
    12:00 p.m. & 5.03 & 5.37 & 0.34 \\
    15:00 (3 p.m.) & 5.69 & 6.03 & 0.34 \\
    \bottomrule
    \end{tabular}
\end{table}

\textbf{Conclusión:} La instalación del sistema solar fotovoltaico produce una reducción en las pérdidas técnicas del sistema, mejorando así la eficiencia global de la red de distribución.

\subsection{Pérdidas Consideradas en Energía Generada}

Para el cálculo de la energía producida por el sistema se consideraron las siguientes pérdidas:

\begin{table}[H]
    \centering
    \caption{Desglose de pérdidas en el sistema fotovoltaico}
    \begin{tabular}{lc}
    \toprule
    \textbf{Tipo de Pérdida} & \textbf{Porcentaje} \\
    \midrule
    Pérdidas por temperatura & 3.00\% \\
    Pérdidas por mismatch & 2.00\% \\
    Pérdidas resistivas & 4.00\% \\
    Pérdidas por conversión DC/AC & 1.00\% \\
    Otras pérdidas asociadas & 16.00\% \\
    Pérdidas por sombreado & 0.00\% \\
    \midrule
    \textbf{Total} & \textbf{26.00\%} \\
    \bottomrule
    \end{tabular}
\end{table}
